\chapter{Conclusion}
\label{chap:conclusion}

Ce projet a abouti à \textbf{PrognoSense}, une plateforme logicielle complète de maintenance
prédictive industrielle qui couvre, de bout en bout, la chaîne de surveillance
conditionnelle : de l'acquisition du signal jusqu'à l'ordre de travail transmis à la GMAO,
en passant par l'extraction d'indicateurs physiquement interprétables, le diagnostic
spectral normalisé ISO~10816, la détection d'anomalies non supervisée, la classification
fine des défauts et le pronostic de durée de vie restante.

Sur le plan scientifique, le travail s'appuie sur une démarche rigoureuse : cinq jeux de
données de référence couvrant un large spectre de défaillances (roulements, engrenages,
balourd, désalignement, jeu mécanique, dégradation de turbomoteur), un benchmark méthodique
de cinq modèles avec une évaluation \emph{sans fuite de données}, et des résultats
significatifs --- 100\,\% sur VBL, 97{,}6\,\% sur CWRU (10 classes), 89{,}6\,\% sur
Mechanical Faults, et un pronostic RUL à 9{,}61 cycles d'erreur moyenne ($R^2 = 0{,}90$) sur
CMAPSS. L'analyse a montré qu'\emph{aucun modèle n'est universellement supérieur}, justifiant
la stratégie de sélection par dataset offerte par la plateforme.

Sur le plan de l'ingénierie, PrognoSense se distingue par sa couche d'\textbf{industrialisation}
: connectivité terrain ouverte (OPC-UA, MQTT), apprentissage du comportement sain propre à
chaque machine, indicateurs de fiabilité \emph{mesurés} (MTBF, MTTR, disponibilité, ROI),
boucle fermée vers la GMAO, et une couche de \textbf{confiance} (explicabilité, calibration,
détection de dérive, journal d'audit, versioning avec retour arrière, mesure réelle des
fausses alarmes). Le copilot conversationnel, ancré sur une base de connaissances normative,
rend cette puissance accessible en langage naturel.

C'est précisément cette \textbf{intégration verticale} --- rigueur algorithmique,
crédibilité normative, actionnabilité opérationnelle et gouvernance de confiance, le tout
dans une architecture ouverte et scalable --- qui distingue PrognoSense des travaux centrés
sur un seul algorithme ou un seul cas d'usage. Le projet ne se contente pas de
\emph{prédire} : il dit \emph{quoi} tombe en panne, \emph{pourquoi} le système le pense,
\emph{depuis quand} il le sait, et \emph{quelle action} entreprendre --- tout en restant
vérifiable.

Les limites identifiées --- validation sur données de terrain, transférabilité du pronostic,
durcissement sécurité --- tracent une feuille de route claire vers une solution
commercialisable. Les perspectives (modèles physics-informed, apprentissage fédéré,
architecture OSA-CBM formelle, écosystème ouvert) confirment le potentiel de la plateforme à
devenir une alternative crédible, ouverte et explicable aux suites propriétaires du marché.

Au-delà du livrable technique, ce projet aura été l'occasion de mettre en œuvre, sur un cas
industriel réaliste et exigeant, l'ensemble de la chaîne de valeur de la science des données
appliquée à la fiabilité : du traitement du signal à l'apprentissage automatique, de
l'architecture logicielle à l'expérience utilisateur, et de la performance brute à la
confiance et à la traçabilité indispensables en milieu industriel.
